\documentclass{article}
\usepackage{apacite}
\usepackage{hyperref}
\usepackage{multicol}
\usepackage{natbib}
\usepackage{sophia}
\usepackage[utf8]{inputenc}
\usepackage{array}
\setlength{\parindent}{0pt}
\usepackage{hyperref}

\usetikzlibrary{calc,patterns,angles,quotes}
\usepackage{pgfplots}
\pgfplotsset{width=15cm,height=10cm,compat=1.9}

\usepackage[english]{babel}
\usepackage{fancyhdr}
\usepackage{lastpage}
\pagestyle{fancy}
\fancyhf{}


\pagestyle{empty}% for cropping
\usepackage{tikz}
\usetikzlibrary{calc}
\usepackage{zref-savepos}
\usepackage{tabu}

\newcounter{NoTableEntry}
\renewcommand*{\theNoTableEntry}{NTE-\the\value{NoTableEntry}}

\newcommand*{\strike}[2]{%
  \multicolumn{1}{#1}{%
    \stepcounter{NoTableEntry}%
    \vadjust pre{\zsavepos{\theNoTableEntry t}}% top
    \vadjust{\zsavepos{\theNoTableEntry b}}% bottom
    \zsavepos{\theNoTableEntry l}% left
    \hspace{0pt plus 1filll}%
    #2% content
    \hspace{0pt plus 1filll}%
    \zsavepos{\theNoTableEntry r}% right
    \tikz[overlay]{%
      \draw
        let
          \n{llx}={\zposx{\theNoTableEntry l}sp-\zposx{\theNoTableEntry r}sp-\tabcolsep},
          \n{urx}={\tabcolsep},
          \n{lly}={\zposy{\theNoTableEntry b}sp-\zposy{\theNoTableEntry r}sp},
          \n{ury}={\zposy{\theNoTableEntry t}sp-\zposy{\theNoTableEntry r}sp}
        in
        (\n{llx}, \n{lly}) -- (\n{urx}, \n{ury})
      ;
    }% 
  }%
}

\usepackage{tikz}
\usetikzlibrary{calc}

\tikzset{
    right angle quadrant/.code={
        \pgfmathsetmacro\quadranta{{1,1,-1,-1}[#1-1]}     % Arrays for selecting quadrant
        \pgfmathsetmacro\quadrantb{{1,-1,-1,1}[#1-1]}},
    right angle quadrant=1, % Make sure it is set, even if not called explicitly
    right angle length/.code={\def\rightanglelength{#1}},   % Length of symbol
    right angle length=2ex, % Make sure it is set...
    right angle symbol/.style n args={3}{
        insert path={
            let \p0 = ($(#1)!(#3)!(#2)$) in     % Intersection
                let \p1 = ($(\p0)!\quadranta*\rightanglelength!(#3)$), % Point on base line
                \p2 = ($(\p0)!\quadrantb*\rightanglelength!(#2)$) in % Point on perpendicular line
                let \p3 = ($(\p1)+(\p2)-(\p0)$) in  % Corner point of symbol
            (\p1) -- (\p3) -- (\p2)
        }
    }
}
\pagestyle{fancy}
\fancyhf{}
\title{Andi's Solution To IYMC Qualification Round}
\author{Jingzhi Andi Xie}
\date{\today}
\lhead{Jingzhi Andi Xie}
\rhead{\today}
\pagestyle{fancy}

 
\cfoot{Page \thepage \hspace{1pt} of \pageref{LastPage}}


\begin{document}

\maketitle


\begin{custom-big}[Problem A]
\[\left(\log\left(3^x\right)-2\log\left(3\right)\right)\cdot\left(x^2-1\right)=0\]
Case 1:
\begin{align*}
    \log\left(3^x\right)-2\log\left(3\right)&=0\\
    \log\left(\frac{3^x}{9}\right)&=0\\
    \frac{3^x}{9}&=1\\
    3^x&=9\\
    x&=2
\end{align*}
Case 2:
\begin{align*}
    x^2-1&=0\\
    \left(x+1\right)\left(x-1\right)&=0\\
    x&=\pm 1
\end{align*}
$\therefore \boxed{x=\pm 1,2}$
\end{custom-big}

\begin{custom-big}[Problem B]
\begin{align*}
    &1+\frac{3}{\pi}+\frac{3}{\pi^2}+\frac{3}{\pi^3}+\cdots\\
    =\ &1+\frac{\frac{3}{\pi}}{1-\frac{1}{\pi}}\\
    =\ &\boxed{1+\frac{3}{\pi-1}}
\end{align*}
\end{custom-big}


\begin{custom-big}[Problem C]
\begin{align*}
    &\log\left[\log\left(3\right)\cdot\left(\log\left(2\right)\cdot\left(\frac{\sqrt{3}-2\sin\left(\pi/3\right)}{\pi^3+1}+1\right)\right)-\log\left(2\right)\log\left(3\right)+\left(-1\right)^{100}\right]\\
    =\ &\log\left[\log\left(3\right)\cdot\left(\log\left(2\right)\cdot\left(\frac{\sqrt{3}-2\left(\frac{\sqrt{3}}{2}\right)}{\pi^3+1}+1\right)\right)-\log\left(2\right)\log\left(3\right)+\left(-1\right)^{100}\right]\\
    =\ &\log\left[\log\left(3\right)\log\left(2\right)-\log\left(2\right)\log\left(3\right)+\left(-1\right)^{100}\right]\\
    =\ &\log\left(1\right)\\
    =\ &\boxed{0}
\end{align*}
\end{custom-big}

\begin{custom-big}[Problem D]
\begin{align*}
    \sigma\left(p^2\right)&=3\ \forall \ p\in\mathbb{P}\left(1,p,p^2\right)\\
\because p&\ge 2\\
\therefore n&\ge 4\\
\therefore \sigma\left(n\right)&<2n\text{ for all prime }p
\end{align*}
\end{custom-big}


\begin{custom-big}[Problem E]
\begin{align*}
    \because \text{Base of the grey}\triangle &=\frac 1 2\text{ Base of the big}\triangle\\
    \therefore\text{From similarity, height of the grey}\triangle &=\frac 1 2\text{ Height of the big}\triangle\\
    \therefore \left[\text{Grey}\triangle\right]&=\frac 14\left[\text{Big}\triangle\right]\\
    \left[\text{Grey}\triangle\right]&=\frac 14\left(\frac{\sqrt{3}}{4}a^2\right)\\
    \left[\text{Grey}\triangle\right]&=\boxed{\frac{\sqrt{3}}{16}a^2}
\end{align*}
\end{custom-big}
\end{document}
